\begin{abstract}
\noindent Survey incomes are reported in round numbers, and tests of unimodality applied to such reports reject for every unimodal income distribution we simulate. We study what the reports identify. Conditional on the rounding level $\ell$, the greatest common divisor $G$ of a report with the least common multiple $W$ of the candidate levels satisfies $\Prob(G=d\mid\ell)=\ell\,\varphi(W/d)/W$, with $\varphi$ Euler's totient. The rounding probabilities are therefore identified under a regularity condition on the residues of income modulo $W$, and the residue of a report modulo $W$ is uniformly distributed given its divisibility pattern, which yields a specification test free of nuisance parameters. At any width that is a common multiple of the active levels the cell probabilities of latent income are identified up to a second-order boundary term, and binning preserves unimodality; we test unimodality of the cell vector with a moment-selection bootstrap whose asymptotic level is at most nominal at every point of the null, including its boundary. Finer resolution is identified by inverting the reporting map, but the inversion needs the rounding profile conditional on latent income, which the reports identify only coarsely. On American data the specification test rejects the mechanism for every earnings variable, and a thinning in the disposable income of one-person households turns out to be rounding in gross earnings carried through the tax computation.
\end{abstract}

\noindent\textbf{Keywords:} heaping; rounding; modality; excess mass; shape-constrained inference; moment selection.
\medskip

\noindent\textbf{JEL codes:} C12, C14, C81, D31.

